\section{Local factors and affine correlations}
\label{sec:affine-corollaries}

We now deduce correlations on arbitrary fixed progressions and affine
forms from \Cref{thm:elliott}. The main point is to handle the primes
dividing a progression modulus or a dilation without assuming complete
multiplicativity. A finite expansion at these primes will suffice.

\subsection{Stability and finite local expansions}

We first extend \Cref{lem:finite-prime-stability} to character twists
and complex conjugation.

\begin{lemma}[Character twists and conjugation]
\label{lem:affine-stability}
Let $f,G\colon\N\to\mathbb D$ be multiplicative, let $\psi$ be a
fixed Dirichlet character, and let $S$ be a finite set of primes.
If $G(p)=f(p)\psi(p)$ for every $p\notin S$, then for every Dirichlet
character $\chi$, every real $t$, and every $X\ge2$,
\begin{equation}
 \left|D(G,\chi n^{it};X)^2
       -D(f,(\chi\overline\psi)n^{it};X)^2\right|
 \le 2\sum_{p\in S}\frac1p.
 \label{eq:affine-stability}
\end{equation}
Thus uniform nonpretentiousness of $f$ implies that of $G$.
Complex conjugation also preserves uniform nonpretentiousness.
\end{lemma}

\begin{proof}
The product $\chi\overline\psi$ is a Dirichlet character modulo the
least common multiple of the two moduli. At every prime outside $S$,
\[
 G(p)\overline{\chi(p)p^{it}}
 =f(p)\overline{(\chi\overline\psi)(p)p^{it}}.
\]
This identity holds also at primes where a character vanishes. Each
prime in $S$ contributes at most $2/p$ to the difference of squared
distances, proving \eqref{eq:affine-stability}. Take the infimum over
exactly $|t|\le X$ on both sides of the resulting lower bound. For each
fixed $\chi$, the character $\chi\overline\psi$ is fixed, and the
bounded error cannot prevent divergence to infinity.
For conjugation the exact identity is
\[
 D(\overline f,\chi n^{it};X)^2
 =D(f,\overline\chi n^{-it};X)^2.
\]
The interval $|t|\le X$ is unchanged by $t\mapsto-t$.
\end{proof}

A multiplicative function is either identically zero or has value one
at $1$: apply multiplicativity to $(1,n)$. Zero factors give zero
correlations, so in the following local construction we may work with
normalized functions. Write $v_p(m)$ for the exponent of $p$ in $m$.

\begin{lemma}[A finite expansion for dilation]
\label{lem:affine-dilation}
Let $a\ge1$ be an integer and let $f\colon\N\to\mathbb D$ be
multiplicative with $f(1)=1$. Define
\begin{equation}
 F_{a,f}(m)=
 \begin{cases}
 f(m/a),&a\mid m,\\
 0,&a\nmid m.
 \end{cases}
 \label{eq:affine-dilation-function}
\end{equation}
Then $F_{a,f}$ is a signed sum of $2^{\omega(a)}$ normalized,
$1$-bounded multiplicative functions, each agreeing with $f$ at every
prime power whose prime does not divide $a$. If $f$ is uniformly
nonpretentious, every function in this sum is uniformly nonpretentious.
\end{lemma}

\begin{proof}
Let $P(a)$ be the set of prime divisors of $a$. For $p\in P(a)$ set
$\alpha_p=v_p(a)$ and define
\begin{equation}
 A_p(k)=
 \begin{cases}
 0,&0\le k<\alpha_p,\\
 f(p^{k-\alpha_p}),&k\ge\alpha_p.
 \end{cases}
 \label{eq:affine-local-sequence}
\end{equation}
Factorization into pairwise coprime prime powers gives the exact identity
\begin{equation}
 F_{a,f}(m)=
 \prod_{p\in P(a)}A_p(v_p(m))
 \prod_{\substack{p\mid m\\p\notin P(a)}}f(p^{v_p(m)}).
 \label{eq:affine-local-product}
\end{equation}
No identity between $f(p^k)$ and $f(p)^k$ is used.

The obstruction to multiplicativity in this product is that
$A_p(0)=0$. Resolve it by setting
\begin{equation}
 B_p(0)=C_p(0)=1,\qquad
 B_p(k)=A_p(k),\quad C_p(k)=0\quad(k\ge1).
 \label{eq:affine-b-minus-c}
\end{equation}
Then $A_p=B_p-C_p$ at every nonnegative exponent. For
$E\subseteq P(a)$ put
\[
 f_{a,E}(m)=
 \prod_{p\in E}C_p(v_p(m))
 \prod_{p\in P(a)\setminus E}B_p(v_p(m))
 \prod_{\substack{p\mid m\\p\notin P(a)}}f(p^{v_p(m)}).
\]
All local factors have modulus at most one and value one at exponent
zero. If $u$ and $v$ are coprime, at most one of $v_p(u),v_p(v)$ is
positive at each prime; hence $f_{a,E}(uv)=f_{a,E}(u)f_{a,E}(v)$.
Expanding the finite product in \eqref{eq:affine-local-product} yields
\begin{equation}
 F_{a,f}(m)=
 \sum_{E\subseteq P(a)}(-1)^{|E|}f_{a,E}(m).
 \label{eq:affine-dilation-expansion}
\end{equation}
Each summand agrees with $f$ at prime powers outside $P(a)$, so
\Cref{lem:finite-prime-stability} proves the last assertion.
\end{proof}

For $a=1$ the expansion has just one term, namely $f$. For $a>1$
and $m=1$, its signed sum is
$\sum_{E\subseteq P(a)}(-1)^{|E|}=0$, as required. Thus the expansion
also handles the value at $1$, despite normalizing every summand there.

\begin{lemma}[Every residue class]
\label{lem:affine-residue}
Let $l\ge1$ and $0\le b<l$ be integers, and set
\[
 d=\gcd(b,l),\qquad q=l/d,\qquad c=b/d.
\]
For every positive integer $m$,
\begin{equation}
 \1_{m\equiv b\pmod l}
 =\frac1{\varphi(q)}\sum_{\chi\bmod q}
       \overline{\chi(c)}F_{d,\chi}(m).
 \label{eq:affine-character-expansion}
\end{equation}
Each $F_{d,\chi}$ is a signed sum of $2^{\omega(d)}$ normalized,
$1$-bounded multiplicative functions agreeing with $\chi$ at all prime
powers outside $d$. For $q=1$, the character in the formula is the
constant function one, including at the residue $c=0$, and
$\varphi(1)=1$.
\end{lemma}

\begin{proof}
If $d\nmid m$, both sides vanish. If $d\mid m$, the congruence on
the left is equivalent to $m/d\equiv c\pmod q$. When $q>1$, the
residue $c$ is a unit modulo $q$, so character orthogonality on
$(\Z/q\Z)^\times$ gives
\[
 \1_{m/d\equiv c\pmod q}
 =\frac1{\varphi(q)}\sum_{\chi\bmod q}
       \overline{\chi(c)}\chi(m/d).
\]
For a nonunit $m/d$ every character on the right vanishes. The case
$q=1$ is the trivial identity with the stated convention. The finite
expansion now follows from \Cref{lem:affine-dilation}.
\end{proof}

In particular, if $p\mid d$ and $\alpha=v_p(d)$, the local sequence
in this application is
\[
 A_p(k)=0\quad(k<\alpha),\qquad
 A_p(k)=\chi(p)^{k-\alpha}\quad(k\ge\alpha),
 \qquad 0^0=1.
\]
If $p$ also divides $q$, then $\chi(p)=0$ and this sequence is one
at $k=\alpha$ and zero elsewhere. No character value is divided out.
When $b=0$ we have $d=l$ and $q=1$, so the same formula expands
$\1_{l\mid m}$; it includes $l=1$ and $m=1$.

\subsection{Progressions and the affine deduction}

The two expansions just proved reduce the desired averages to finitely
many applications of \Cref{thm:elliott}. We first incorporate a
progression restriction, placing its weight at the argument of a
nonpretentious factor.

\begin{proposition}[Correlations on fixed progressions]
\label{prop:affine-progression}
Let $f_1,f_2\colon\N\to\mathbb D$ be multiplicative, with at least
one uniformly nonpretentious. For fixed distinct nonnegative integers
$h_1,h_2$, an integer $l\ge1$, and any integer $b$,
\begin{equation}
 \frac1X\sum_{1\le n\le X}
 f_1(n+h_1)f_2(n+h_2)\1_{n\equiv b\pmod l}
 \longrightarrow0
 \qquad(X\longrightarrow\infty).
 \label{eq:affine-progression-limit}
\end{equation}
The limit holds through all real $X$, including nonunit residue classes.
\end{proposition}

\begin{proof}
We may assume both functions are nonzero. Choose $j\in\{1,2\}$ for
which $f_j$ is uniformly nonpretentious. Since
\[
 \1_{n\equiv b\pmod l}
 =\1_{n+h_j\equiv b+h_j\pmod l},
\]
\Cref{lem:affine-residue}, using the least nonnegative representative
of $b+h_j$ modulo $l$, expands the indicator as a function of
$m=n+h_j$ into finitely many bounded multiplicative functions, each
agreeing with a fixed Dirichlet character outside a fixed finite prime
set. Their products with $f_j$ are uniformly nonpretentious by
\Cref{lem:affine-stability}. Apply \Cref{thm:elliott} to each product
at shift $h_j$ and the unchanged other factor at its distinct shift,
and sum the finitely many conclusions. This proves
\eqref{eq:affine-progression-limit} at integer $X$.
For real $X$, replacing the normalization $X$ by $\lfloor X\rfloor$
changes the average by at most $1/X$, proving the assertion.
\end{proof}

\begin{proof}[Proof of \Cref{cor:affine}]
Zero factors again give an immediate conclusion. Otherwise set
\begin{equation}
 \Delta=a_1b_2-a_2b_1,\qquad
 l=a_1a_2,\qquad
 c_0=\min(a_2b_1,a_1b_2),\qquad h=|\Delta|\ge1.
 \label{eq:affine-parameters}
\end{equation}
For $\Delta>0$ put
$U=F_{a_2,f_1}$ and $V=F_{a_1,f_2}$; for $\Delta<0$ put
$U=F_{a_1,f_2}$ and $V=F_{a_2,f_1}$. In either case,
\begin{equation}
 U(ln+c_0)V(ln+c_0+h)
 =f_1(a_1n+b_1)f_2(a_2n+b_2).
 \label{eq:affine-exact-identity}
\end{equation}
Indeed, when $\Delta>0$ the two arguments on the left are
$a_2(a_1n+b_1)$ and $a_1(a_2n+b_2)$; when $\Delta<0$ their order
is reversed.

Expand $U$ and $V$ by \Cref{lem:affine-dilation}. Each pair of
multiplicative components has a uniformly nonpretentious factor,
inherited from the corresponding original function. Applying
\Cref{prop:affine-progression} termwise with shifts $0,h$ gives
\begin{equation}
 \frac1X\sum_{1\le m\le X}
 U(m)V(m+h)\1_{m\equiv c_0\pmod l}\longrightarrow0.
 \label{eq:affine-expanded-limit}
\end{equation}
Put $T(m)=U(m)V(m+h)\1_{m\equiv c_0\pmod l}$ and $X=lN+c_0$.
The integers in this class satisfying $c_0<m\le X$ are precisely
$ln+c_0$, $1\le n\le N$. Therefore
\begin{align}
 \frac1N\sum_{n=1}^N f_1(a_1n+b_1)f_2(a_2n+b_2)
 &=\frac XN\left(\frac1X\sum_{1\le m\le X}T(m)\right)
   -\frac1N\sum_{1\le m\le c_0}T(m).
 \label{eq:affine-endpoints}
\end{align}
The first term tends to zero because $X/N\to l$ and
\eqref{eq:affine-expanded-limit} holds. The second has modulus at most
$c_0/N$ and is empty when $c_0=0$. This proves the corollary, including
zero intercepts and coefficients with common factors.
\end{proof}

\subsection{Liouville and M\"obius correlations}

To apply the preceding conclusions to the Liouville function, we verify
the nonpretentiousness hypothesis over its full growing height interval.
The required uniform estimate is supplied by
Matom\"aki, Radziwi{\l}{\l}, and Tao's version of an argument of
Granville and Soundararajan \cite[Appendix~C]{MRT15}.

\begin{lemma}[Uniform nonpretentiousness of Liouville]
\label{lem:affine-liouville-unp}
For every fixed Dirichlet character $\chi$,
\begin{equation}
 \inf_{|t|\le X}D(\lambda,\chi n^{it};X)
 \ge \frac14\sqrt{\log\log X}-O_\chi(1)
 \qquad(X\ge100).
 \label{eq:affine-liouville-unp}
\end{equation}
In particular, $\lambda$ is uniformly nonpretentious.
\end{lemma}

\begin{proof}
For a real $1$-bounded multiplicative function $f$,
\cite[Lemma~C.1]{MRT15} gives
\begin{equation}
 D(f,\chi n^{it};X)
 \ge\frac14\sqrt{\log\log X}-O_\chi(1)
 \quad(1\le|t|\le X).
 \label{eq:affine-real-large-height}
\end{equation}
When $\chi^2$ is nonprincipal the same bound holds throughout
$|t|\le X$. Thus taking $f=\lambda$ proves the assertion for nonreal
$\chi$ and handles $1\le|t|\le X$ for all $\chi$.
For real $\chi$, the other part of the cited lemma gives
\begin{equation}
 D(\lambda,\chi n^{it};X)
 \ge\frac13D(\lambda,\chi;X)-O(1)
 \quad(|t|\le1).
 \label{eq:affine-real-small-height}
\end{equation}
It remains to estimate the distance at $t=0$ for these real characters.

If $\chi$ is nonprincipal, the prime number theorem in arithmetic
progressions for its fixed modulus yields, for some $c_\chi>0$,
\[
 A_\chi(u):=\sum_{p\le u}\chi(p)\log p
 =O_\chi\bigl(u\exp(-c_\chi\sqrt{\log u})\bigr).
\]
Partial summation gives
\[
 \sum_{p\le X}\frac{\chi(p)}p
 =\frac{A_\chi(X)}{X\log X}
   +\int_2^X A_\chi(u)
     \frac{\log u+1}{u^2(\log u)^2}\,du
 =O_\chi(1),
\]
since the integral converges absolutely as $X\to\infty$.
As $\lambda(p)=-1$, Mertens' estimate now gives
\[
 D(\lambda,\chi;X)^2
 =\sum_{p\le X}\frac{1+\chi(p)}p
 =\log\log X+O_\chi(1).
\]
If instead $\chi$ is principal modulo $q$, then
\[
 D(\lambda,\chi;X)^2
 =2\sum_{p\le X}\frac1p
  -\sum_{\substack{p\mid q\\p\le X}}\frac1p
 =2\log\log X+O_q(1).
\]
In both cases \eqref{eq:affine-real-small-height} is at least
$\frac13\sqrt{\log\log X}-O_\chi(1)$. Combining this with
\eqref{eq:affine-real-large-height} proves
\eqref{eq:affine-liouville-unp} on the entire interval $|t|\le X$.
\end{proof}

Thus $\lambda$ satisfies the exact hypothesis of the qualitative theorem.
Applying \Cref{prop:affine-progression,cor:affine} with
$f_1=f_2=\lambda$ independently recovers ordinary cancellation in
every fixed residue class and along every fixed nonproportional affine
pair. \Cref{q:progression,thm:q-affine} establish these conclusions
with an absolute power-of-logarithm saving.

The qualitative theorem also applies to the M\"obius function, defined by
\[
 \mu(n)=
 \begin{cases}
 \lambda(n),&n\text{ is squarefree},\\
 0,&n\text{ is not squarefree}.
 \end{cases}
\]
This function is multiplicative and has $\mu(p)=\lambda(p)=-1$ at
every prime. Since the distance uses only prime values,
\[
 D(\mu,\chi n^{it};X)=D(\lambda,\chi n^{it};X)
\]
for every Dirichlet character $\chi$, real $t$, and $X\ge2$.
The preceding lemma therefore proves uniform nonpretentiousness of
$\mu$ with the same growing twist range. The conclusions of
\Cref{prop:affine-progression,cor:affine} hold for any choice
$f_1,f_2\in\{\lambda,\mu\}$. This proves ordinary cancellation for two-point M\"obius
and mixed M\"obius--Liouville correlations in every fixed
residue class and along every fixed nonproportional affine pair.
For products containing a M\"obius factor, this deduction gives
convergence without a quantitative rate.
